% GENERATED FILE. Edit canonical lessons, then run npm run generate:books.
% canonical-source-hash: fnv1a64:4587ebee4ab673a6
% canonical-lessons: PA-C106-kite, PA-C106-kade, PA-C106-harek, PA-C106-ajiha, PA-C106-har2

\chapter{Somewhere, Ever, Each, Such, Every}
\label{ch:pa-somewhere-ever-each-such-every}
\hlchaptermodality{106}

\begin{chapteropening}
\textbf{By the end of this chapter:} \emph{I can say somewhere, ever, each, such and every.}

Complete the last lesson of chapter 106: i can say somewhere, ever, each, such and every.
\end{chapteropening}

\section[\texorpdfstring{\pa{ਕਿਤੇ} (kite)}{ਕਿਤੇ (kite)}]{\pa{ਕਿਤੇ} (kite) — somewhere}

\label{lesson:PA-C106-kite}



\begin{quote}
Before the new one: say the Punjabi for Magh, then the Punjabi for Phaggan.
\end{quote}

\subsection*{You'll want to know: \pa{ਕਿਤੇ}}
\textbf{\pa{ਕਿਤੇ}} — \emph{kite} — \textquotedblleft{}somewhere\textquotedblright{}.

\begin{grammarlens}[title={What you've built}]
One more word for this chapter.
\end{grammarlens}

\subsection*{Your turn}
\begin{itemize}
\raggedright
  \item \emph{Say it:} \emph{kite}
  \item \emph{Say it:} \emph{kite}, once more
  \item \emph{Say it:} say \emph{kite}
  \item \emph{Recall:} say the Punjabi for Magh, then the Punjabi for Phaggan, then say \emph{kite} again
\end{itemize}

\begin{tcolorbox}[breakable,colback=teal!4,colframe=teal!35!black,title={Before you move on}]
What does \pa{ਕਿਤੇ} mean? (\textquotedblleft{}Somewhere\textquotedblright{}.) Say it once more.
\end{tcolorbox}
\section[\texorpdfstring{\pa{ਕਦੇ} (kade)}{ਕਦੇ (kade)}]{\pa{ਕਦੇ} (kade) — ever, sometime}

\label{lesson:PA-C106-kade}



\begin{quote}
Before the new one: say the Punjabi for Phaggan, then the Punjabi for somewhere.
\end{quote}

\subsection*{You'll want to know: \pa{ਕਦੇ}}
\textbf{\pa{ਕਦੇ}} — \emph{kade} — \textquotedblleft{}ever, sometime\textquotedblright{}.

\begin{grammarlens}[title={What you've built}]
One more word for this chapter.
\end{grammarlens}

\subsection*{Your turn}
\begin{itemize}
\raggedright
  \item \emph{Say it:} \emph{kade}
  \item \emph{Say it:} \emph{kade}, once more
  \item \emph{Say it:} say \emph{kade}
  \item \emph{Recall:} say the Punjabi for Phaggan, then the Punjabi for somewhere, then say \emph{kade} again
\end{itemize}

\begin{tcolorbox}[breakable,colback=teal!4,colframe=teal!35!black,title={Before you move on}]
What does \pa{ਕਦੇ} mean? (\textquotedblleft{}Ever, sometime\textquotedblright{}.) Say it once more.
\end{tcolorbox}
\section[\texorpdfstring{\pa{ਹਰੇਕ} (harek)}{ਹਰੇਕ (harek)}]{\pa{ਹਰੇਕ} (harek) — each}

\label{lesson:PA-C106-harek}



\begin{quote}
Before the new one: say the Punjabi for somewhere, then the Punjabi for ever.
\end{quote}

\subsection*{You'll want to know: \pa{ਹਰੇਕ}}
\textbf{\pa{ਹਰੇਕ}} — \emph{harek} — \textquotedblleft{}each\textquotedblright{}.

\begin{grammarlens}[title={What you've built}]
One more word for this chapter.
\end{grammarlens}

\subsection*{Your turn}
\begin{itemize}
\raggedright
  \item \emph{Say it:} \emph{harek}
  \item \emph{Say it:} \emph{harek}, once more
  \item \emph{Say it:} say \emph{harek}
  \item \emph{Recall:} say the Punjabi for somewhere, then the Punjabi for ever, then say \emph{harek} again
\end{itemize}

\begin{tcolorbox}[breakable,colback=teal!4,colframe=teal!35!black,title={Before you move on}]
What does \pa{ਹਰੇਕ} mean? (\textquotedblleft{}Each\textquotedblright{}.) Say it once more.
\end{tcolorbox}
\section[\texorpdfstring{\pa{ਅਜਿਹਾ} (ajihā)}{ਅਜਿਹਾ (ajihā)}]{\pa{ਅਜਿਹਾ} (ajihā) — such}

\label{lesson:PA-C106-ajiha}



\begin{quote}
Before the new one: say the Punjabi for ever, then the Punjabi for each.
\end{quote}

\subsection*{You'll want to know: \pa{ਅਜਿਹਾ}}
\textbf{\pa{ਅਜਿਹਾ}} — \emph{ajihā} — \textquotedblleft{}such\textquotedblright{}.

\begin{grammarlens}[title={What you've built}]
One more word for this chapter.
\end{grammarlens}

\subsection*{Your turn}
\begin{itemize}
\raggedright
  \item \emph{Say it:} \emph{ajihā}
  \item \emph{Say it:} \emph{ajihā}, once more
  \item \emph{Say it:} say \emph{ajihā}
  \item \emph{Recall:} say the Punjabi for ever, then the Punjabi for each, then say \emph{ajihā} again
\end{itemize}

\begin{tcolorbox}[breakable,colback=teal!4,colframe=teal!35!black,title={Before you move on}]
What does \pa{ਅਜਿਹਾ} mean? (\textquotedblleft{}Such\textquotedblright{}.) Say it once more.
\end{tcolorbox}
\section[\texorpdfstring{\pa{ਹਰ} (har)}{ਹਰ (har)}]{\pa{ਹਰ} (har) — every}

\label{lesson:PA-C106-har2}



\begin{quote}
Before the new one: say the Punjabi for each, then the Punjabi for such.
\end{quote}

\subsection*{You'll want to know: \pa{ਹਰ}}
\textbf{\pa{ਹਰ}} — \emph{har} — \textquotedblleft{}every\textquotedblright{}.

\begin{grammarlens}[title={What you've built}]
One more word for this chapter.
\end{grammarlens}

\subsection*{Your turn}
\begin{itemize}
\raggedright
  \item \emph{Say it:} \emph{har}
  \item \emph{Say it:} \emph{har}, once more
  \item \emph{Say it:} say \emph{har}
  \item \emph{Recall:} say the Punjabi for each, then the Punjabi for such, then say \emph{har} again
\end{itemize}

\begin{tcolorbox}[breakable,colback=teal!4,colframe=teal!35!black,title={Before you move on}]
What does \pa{ਹਰ} mean? (\textquotedblleft{}Every\textquotedblright{}.) Say it once more.
\end{tcolorbox}
